\section{Empirical Results: the Agglutinative Compounding Effect}
\label{sec:results}

\S\ref{sec:framework} gave the framework two independent predictions. The first (equation~\ref{eq:rebate} in its rank form) asserts that the parameter rebate morphology-aware tokenization delivers over a BPE baseline is monotonically ordered by $\rho(L) \cdot H(L)$. The second (the same equation in its magnitude form) asserts that the proportionality holds with a single coefficient $\alpha$ constant across languages. This section reports a six-model Phase~1 experiment testing both, and extends it to a Chinchilla-ratio three-rung scale ladder (\S\ref{sec:results-ladder}) that probes whether the observed phenomenon survives at properly-trained scales.

\subsection{Design}
\label{sec:results-design}

The Phase~1 mini experiment trained six decoder-only transformer models (English, Arabic, and Turkish, each under both BPE and morphology-aligned tokenization) under identical architecture, data, optimizer, and training budget, per an experimental contract frozen before implementation (\texttt{experimental\_contract.md}). Each model used 4 layers, 128-dimensional embeddings, 4 attention heads, a 128-token context, and was trained for $5 \times 10^6$ tokens with AdamW at fixed hyperparameters \citep{loshchilov2019}, yielding parameter counts in the $2 \times 10^6$ range. Within-language parity is strict: variants differ only in how text becomes token identifiers and, for the morph variant, in an additional feature-bundle embedding summed at the input layer. The configuration sits substantially below the reasoning-emergence regime and is approximately eight times Chinchilla-suboptimal on tokens; these facts bound the interpretability below, and the interpretations in \S\ref{sec:results-interpretations} turn on them.

\FloatBarrier

\subsection{Ordinal validation}
\label{sec:results-ordinal}

Table~\ref{tab:ordinal} reports the test cross-entropy of each model on a held-out sample of the same Wikipedia corpus used for training, together with the framework's predicted ordering variable $\rho(L) \cdot H(L)$ computed in \S\ref{sec:framework}. The ordering by $\rho \cdot H$ (English $1.28 <$ Arabic $3.20 <$ Turkish $4.20$) matches the ordering by observed $\Delta\mathcal{L}$ monotonically. Translated into fractional perplexity reduction, this is the 1\,\% / 14\,\% / 57\,\% gradient that first motivated the investigation. The framework's qualitative prediction therefore survives its test. The three languages span more than an order of magnitude in $\Delta\mathcal{L}$, and the ordering is preserved across the full range. Crucially, $\rho(L) \cdot H(L)$ is computed from the grammar engine alone, without reference to any training run: the match is a forward prediction from structural information, not a back-fit of a chosen metric.

\begin{table}[!htbp]
\centering
\small
\begin{tabular}{lrrrrrrr}
\toprule
Language & $\mathcal{L}_\mathrm{BPE}$ & $\mathcal{L}_\mathrm{morph}$ & $\Delta\mathcal{L}$ & $\Delta\mathcal{L}/\mathcal{L}_\mathrm{BPE}$ & $H(L)$ & $\rho(L)$ & $\rho\cdot H$ \\
\midrule
English & 5.803 & 5.792 & 0.0107 & 0.18\% & 1.43 & 0.89 & \textbf{1.28} \\
Arabic  & 6.413 & 6.263 & 0.1497 & 2.33\% & 5.44 & 0.59 & \textbf{3.20} \\
Turkish & 4.187 & 3.344 & 0.8427 & 20.13\% & 4.97 & 0.85 & \textbf{4.20} \\
\bottomrule
\end{tabular}
\caption{Phase~1 test cross-entropy and the framework's predicted ordering variable $\rho(L) \cdot H(L)$. $\Delta\mathcal{L}$ and $\rho \cdot H$ are monotonically co-ordered across all three languages.}
\label{tab:ordinal}
\end{table}

\begin{figure}[!htbp]
\centering
\includegraphics[width=0.82\linewidth]{fig_01_ordinal_validation.pdf}
\caption{Ordinal validation of the parameter-rebate prediction at Phase~1 scale. Observed $\Delta\mathcal{L}$ (log scale) plotted against the engine-derived $\rho(L) \cdot H(L)$ for the three Phase~1 languages. Dashed line indicates monotonic ordering.}
\label{fig:ordinal}
\end{figure}

\FloatBarrier
\subsection{Magnitude analysis at Phase 1 scale}
\label{sec:results-magnitude}

We next convert each observed $\Delta\mathcal{L}$ into an implied parameter rebate $\Delta N(L)$ and, via equation~\ref{eq:rebate}, into an implied coefficient $\alpha$. The conversion from loss to parameters uses the \citet{hoffmann2022} scaling-law derivative evaluated at the Phase~1 operating point:

\begin{equation*}
\partial\mathcal{L} / \partial N = -\alpha_N \cdot (\mathcal{L} - E) / N, \qquad
\Delta N = \Delta\mathcal{L} \cdot N / [\alpha_N \cdot (\mathcal{L}_\mathrm{BPE} - E)],
\end{equation*}

with $\alpha_N = 0.34$ and $N = 2 \times 10^6$. The irreducible entropy floor $E$ is not directly measurable from six runs; we therefore report results under three values spanning the plausible range: $E = 0$, $E = 1.7$ nats (Chinchilla-era English estimate), and a language-specific $E_L$ pinning $\mathcal{L}_\mathrm{BPE} - E_L = 1$ nat. Under the Chinchilla value, the three languages imply $\alpha_\mathrm{En} \approx 12{,}000$, $\alpha_\mathrm{Ar} \approx 58{,}000$, $\alpha_\mathrm{Tr} \approx 474{,}000$ parameters per bit: a spread of just under 40-fold. Figure~\ref{fig:alpha-dispersion} shows that the pattern survives any $E$ choice, and in every case the implied $\alpha$ grows with $\rho \cdot H$ itself. Assumption~A1 (a single cross-lingual $\alpha$) is not tenable at Phase~1 scale.

\begin{figure}[!htbp]
\centering
\includegraphics[width=0.88\linewidth]{fig_02_alpha_dispersion.pdf}
\caption{Dispersion of implied $\alpha$ across the three Phase~1 languages under three choices of the entropy floor $E$. Implied $\alpha$ spans approximately forty-fold under any $E$ choice, a dispersion the linear Assumption A1 cannot accommodate at Phase~1 scale.}
\label{fig:alpha-dispersion}
\end{figure}

\FloatBarrier
\subsection{Naming the phenomenon}
\label{sec:results-naming}

The consistent direction of the deviation (implied $\alpha$ accelerating with $\rho \cdot H$) is the empirical shape the framework did not anticipate. We name it for reference:

\begin{quote}
\textbf{The Agglutinative Compounding Effect.} Under identical architecture, data, and training budget, the parameter rebate that morphology-aligned tokenization delivers over a BPE baseline grows super-linearly in the effective grammatical information content $\rho(L) \cdot H(L)$ of the language. Languages whose morphology approaches the high-recoverability end of the typological spectrum (agglutinative systems in the classical sense) receive disproportionately large savings per bit of grammatical information they expose.
\end{quote}

The name reflects the typological pole at which the effect is most visible: Turkish shows it most strongly, Arabic (templatic fusion suppresses $\rho$) in attenuated form, English (where $H$ itself is small) least. The claim is not that agglutination \emph{per se} is causal (equation~\ref{eq:rho} defines $\rho$ without reference to typology); agglutinative systems simply maximise $\rho \cdot H$ and so provide the cleanest site of observation. At Phase~1 scale the effect is a three-point observation; naming it asserts it is the thing for which the framework must account, not that it generalises.

\subsection{Two interpretations}
\label{sec:results-interpretations}

Two interpretations of the observed super-linearity are consistent with Phase~1 evidence, and each has very different consequences for the theory and for practice.

\paragraph{Interpretation A: genuine framework non-linearity.}
The assumption A1 of \S\ref{sec:framework-A1} is wrong. The true relationship between grammatical information content and the parameter cost of distributional grammar-learning is not linear in $\rho \cdot H$, but super-linear. The physical basis would be that high-$\rho$ languages permit a cleaner separation of grammatical from semantic learning, and that separation produces compounding efficiencies the linear form fails to capture. Under Interpretation~A, the Agglutinative Compounding Effect is a real property of morphology-aware language modelling.

\paragraph{Interpretation B: scaling-law breakdown at under-trained scale.}
The assumption A1 is fine, but the Hoffmann derivative is not reliable at Phase~1's operating point. Phase~1 trained at approximately $5 \times 10^6$ tokens on $2 \times 10^6$ parameters; Chinchilla-optimal for this parameter count is approximately $4 \times 10^7$ tokens. In the data-starved regime the loss-to-parameter derivative is systematically wrong in a language-dependent way: languages whose BPE baseline sits further from the achievable floor amplify the apparent $\Delta N$ more than languages close to their floor. Under Interpretation~B, the Agglutinative Compounding Effect dissolves at Chinchilla-optimal training and $\alpha$ becomes constant.

Phase~1 cannot choose between A and B. The framework supplies both the question and the discriminating test: a scale ladder that climbs from the Phase~1 regime toward Chinchilla-optimal training. \S\ref{sec:results-ladder} reports the three-rung result.

\FloatBarrier
\subsection{Chinchilla-ratio scale ladder}
\label{sec:results-ladder}

To discriminate between Interpretations A and B, we extend Phase~1 into a Chinchilla-ratio scale ladder. Each rung targets $\approx 20$ training tokens per baseline parameter, placing the loss-to-parameter derivative in the regime where the Hoffmann coefficients were fit. Three rungs are reported: rung~1 ($\approx 620$k params, $5 \times 10^6$ tokens), rung~2 doubling both, and rung~3 doubling again. Rung~4 ($\approx 5$M params, $4 \times 10^7$ tokens) is planned but not included. Morph-regime models carry the same 50k-cap morph vocabulary as Phase~1 and therefore have $\approx 5\times$ more effective parameters than their baseline counterparts, an asymmetry that biases the comparison \emph{against} showing a morph advantage (any observed positive $\Delta\mathcal{L}$ is a conservative lower bound). A parameter-matched follow-up is scoped in \S\ref{sec:conclusion}.

Table~\ref{tab:ladder3} reports the paired results.

\begin{table}[!htbp]
\centering
\small
\begin{tabular}{lrcrrrr}
\toprule
Language & $\rho \cdot H$ & Rung & $\mathcal{L}_\mathrm{BPE}$ & $\mathcal{L}_\mathrm{morph}$ & $\Delta\mathcal{L}$ & $\Delta$-PPL \% \\
\midrule
English & 1.29 & 1 & 6.150 & 6.028 & \textbf{$+0.122$} & $11.5$\,\%  \\
English & 1.29 & 2 & 5.568 & 5.689 & \textbf{$-0.122$} & $-12.9$\,\% \\
English & 1.29 & 3 & 5.007 & 5.345 & \textbf{$-0.337$} & $-40.1$\,\% \\
\midrule
Arabic  & 3.20 & 1 & 6.753 & 6.535 & \textbf{$+0.218$} & $19.6$\,\%  \\
Arabic  & 3.20 & 2 & 6.140 & 6.092 & \textbf{$+0.048$} & $4.7$\,\%   \\
Arabic  & 3.20 & 3 & 5.432 & 5.690 & \textbf{$-0.258$} & $-29.5$\,\% \\
\midrule
Turkish & 4.37 & 1 & 4.624 & 3.853 & \textbf{$+0.770$} & $53.7$\,\%  \\
Turkish & 4.37 & 2 & 3.897 & 3.140 & \textbf{$+0.757$} & $53.1$\,\%  \\
Turkish & 4.37 & 3 & 3.339 & 2.803 & \textbf{$+0.535$} & $41.5$\,\%  \\
\bottomrule
\end{tabular}
\caption{Three-rung Chinchilla-ratio scale ladder, all three Phase~1 languages. English inverts at rung~2; Arabic inverts at rung~3; Turkish remains substantially positive throughout, with $\Delta\mathcal{L}$ decaying from $+0.770$ to $+0.535$ as the baseline trains closer to its entropy floor.}
\label{tab:ladder3}
\end{table}

\begin{figure}[!htbp]
\centering
\includegraphics[width=0.95\linewidth]{fig_06_ordinal_per_rung.pdf}
\caption{Ordinal consistency across the three-rung ladder. Observed $\Delta\mathcal{L}$ plotted against engine-derived $\rho(L) \cdot H(L)$ on a shared linear $y$-axis; the dotted line marks the $\Delta\mathcal{L} = 0$ boundary below which morph underperforms BPE. English crosses below zero at rung~2; Arabic crosses at rung~3; Turkish stays well above at every rung. The engine's ordering $\text{English} < \text{Arabic} < \text{Turkish}$ in $\rho \cdot H$ is preserved at every rung, including the sign-flipped regimes.}
\label{fig:ladder-ordinal}
\end{figure}

\paragraph{Ordinal prediction: robust at every rung.}
The ordering $\Delta\mathcal{L}(\text{English}) < \Delta\mathcal{L}(\text{Arabic}) < \Delta\mathcal{L}(\text{Turkish})$ is preserved at rung~1, rung~2, and rung~3, matching the engine-derived $\rho \cdot H$ ordering ($1.29 < 3.20 < 4.37$) without violation at any scale. The framework's qualitative prediction survives the full three-rung extension cleanly, including into the sign-flipped regime where English and Arabic cross below zero: the ordering still holds, with $-0.337 < -0.258 < +0.535$.

\begin{figure}[!htbp]
\centering
\includegraphics[width=0.80\linewidth]{fig_05_alpha_vs_scale.pdf}
\caption{Implied per-bit coefficient $\alpha$ across the three-rung ladder, symlog $y$-axis. Positive-$\alpha$ rungs sit in the upper half; rungs where $\Delta\mathcal{L}$ has inverted produce negative implied $\alpha$ and sit in the lower half (ringed in red). Turkish's $\alpha$ rises across scales; Arabic crashes below zero by rung~3; English crashes by rung~2. A single cross-lingual constant $\alpha$, as Assumption~A1 posits, would have shown a horizontal band; the observed trajectories are typology-dependent.}
\label{fig:alpha-vs-scale}
\end{figure}

\paragraph{Magnitude: typology-dependent decay rates.}
All three languages' $\Delta\mathcal{L}$ compresses with scale, but at dramatically different rates. English collapses fastest and inverts by rung~2 ($+0.12 \to -0.12 \to -0.34$). Arabic holds through rung~2 but inverts at rung~3 ($+0.22 \to +0.05 \to -0.26$). Turkish compresses too, but only from $+0.770$ at rung~1 to $+0.535$ at rung~3, staying well above zero at rung~3 where English has collapsed to $-0.34$ and Arabic to $-0.26$. The compression is ordered by $\rho \cdot H$ in reverse: low-$\rho \cdot H$ languages' rebates fail first, high-$\rho \cdot H$ languages' rebates last longest.

\paragraph{The Turkish rung-3 result is the decisive datum.}
The morph-vocab-cap asymmetry (50k morph vocabulary versus baseline's 8k BPE vocabulary, producing a $\sim 4.6\times$ parameter inflation on the morph leg) applies identically across all three languages. It alone makes the sign flips in English and Arabic at the later rungs hard to attribute cleanly to the morph advantage disappearing (as Interpretation~B would predict) versus the confound overwhelming it. But Turkish's morph leg sits under exactly the same asymmetric inflation, and Turkish still delivers $+0.535$ nats at rung~3. For Turkish, then, the true morph-mechanism benefit is large enough to absorb the confound and still produce a substantive positive $\Delta\mathcal{L}$; for Arabic and English, it is not. This is strong evidence that the morph advantage at the high-$\rho \cdot H$ pole is a real, mechanism-driven phenomenon of the language--tokenizer pairing (the content Interpretation~A predicts), even as its magnitude decays with scale.

\paragraph{Arabic and English: Interpretation B plus vocab-cap confound, indistinguishable on this data.}
The inversions at rung~2 onward for English and at rung~3 for Arabic are consistent with two mutually-exclusive readings: either the morph advantage for lower-$\rho \cdot H$ languages was genuinely a Chinchilla-suboptimal-baseline artefact that vanishes at well-trained scales (Interpretation~B), or the vocab-cap confound dominates once the true morph advantage is small enough not to overshadow it. The current experimental design cannot separate the two for these languages. A parameter-matched ladder (morph vocabulary capped to baseline's 8k, or morph embedding dimension reduced to compensate) is the discriminating test, scoped as follow-up in \S\ref{sec:conclusion}.

\paragraph{Consequences for the Agglutinative Compounding Effect.}
The Effect is best read as a typology-dependent decay rather than a starvation artefact or a scale-invariant phenomenon. At Phase-1 scale the per-bit rebate appears super-linear because the compression rate toward convergence is inversely proportional to $\rho \cdot H$; at well-trained scales the same hierarchy produces the inversion pattern visible in Table~\ref{tab:ladder3}, with the ordinal property of $\rho \cdot H$ preserved throughout. The framework-level finding is therefore that the morphological-typology axis predicts both cross-linguistic ordering \emph{and} the rate at which each language's rebate compresses as its baseline approaches its entropy floor.

\FloatBarrier
\subsection{The Templatic Lexical Surplus: a two-tier Arabic experiment}
\label{sec:results-tier2}

\S\ref{sec:discussion-limitations} flags a specific gap in the framework: the structural recoverability coefficient $\rho(L)$ measures only what a surface-only parser can decode, but Arabic's templatic morphology stores a substantial fraction of its grammatical information behind lexical dependencies, accessible only with root-lexicon access. To test whether this hidden axis is real and quantifiable, we ran a side experiment: a 2nd-tier Arabic morph variant that adds an explicit pure-root embedding stream on top of the 1st-tier surface-bundle and feature-bundle streams.

The 2nd-tier architecture sums three embeddings at the input layer:
\[
x = \mathrm{tok\_emb}[\text{surface\_bundle}] + \mathrm{feat\_emb}[\text{bundle}] + \mathrm{root\_emb}[\text{root}].
\]
The root embedding is trained from scratch on a 20{,}000-entry vocabulary of distinct triconsonantal roots extracted by the Arabic grammar engine. Words sharing a root (the \emph{k-t-b} family, the \emph{s-l-m} family) receive a tied signal the 1st-tier model cannot access directly.

\begin{table}[!htbp]
\centering
\small
\begin{tabular}{crrrrrr}
\toprule
Rung & $\mathcal{L}_\mathrm{BPE}$ & $\mathcal{L}_\mathrm{1st\text{-}tier}$ & $\mathcal{L}_\mathrm{2nd\text{-}tier}$ & $\Delta\mathcal{L}_\mathrm{1st}$ & $\Delta\mathcal{L}_\mathrm{surplus}$ & Cumulative $\Delta\mathcal{L}$ \\
 & & & & (1st -- BPE) & (2nd -- 1st) & (2nd -- BPE) \\
\midrule
1 & 6.753 & 6.535 & 6.453 & $+0.218$ & \textbf{$+0.082$} & $+0.300$ \\
2 & 6.140 & 6.092 & 6.015 & $+0.048$ & \textbf{$+0.077$} & $+0.125$ \\
3 & 5.432 & 5.690 & 5.632 & $-0.258$ & \textbf{$+0.058$} & $-0.200$ \\
\bottomrule
\end{tabular}
\caption{Arabic three-rung results across BPE baseline, 1st-tier morph, and 2nd-tier morph (root-stream extension). The 1st-tier rebate over baseline collapses and inverts by rung~3, consistent with the broader ladder pattern in \S\ref{sec:results-ladder}. The 2nd-tier \emph{surplus} over 1st-tier (the additional rebate attributable specifically to the explicit root-identity stream) decays only mildly ($+0.082 \to +0.077 \to +0.058$) and remains clearly positive at every rung.}
\label{tab:tier2-ar}
\end{table}

At rung~1 the explicit root stream reduces Arabic's test cross-entropy by an additional 0.082 nats beyond 1st-tier morph, roughly 38\,\% more than 1st-tier alone delivered over baseline. The surplus is nearly preserved at rung~2 ($+0.077$) and decays only slightly at rung~3 ($+0.058$), where the 1st-tier rebate has already inverted to $-0.258$. The 2nd-tier surplus is therefore more scale-robust than the 1st-tier rebate it supplements, and is the only component of Arabic's total rebate over BPE that survives the confound-driven inversion at rung~3. We name this additional rebate the \textbf{Templatic Lexical Surplus}:

\begin{quote}
\textbf{The Templatic Lexical Surplus.} For templatic languages, a portion of the per-word grammatical information content is stored behind lexical dependencies rather than on the surface. Measured as the additional test-loss reduction $\Delta\mathcal{L}_{\mathrm{surplus}}$ delivered by layering an explicit root-identity embedding on top of 1st-tier morph tokenisation, the Templatic Lexical Surplus quantifies the component of $\rho \cdot H$ that the surface-only $\rho$ filter of \S\ref{sec:framework-rho} rejects. Across the three-rung Chinchilla-ratio scale ladder we find $\Delta\mathcal{L}_{\mathrm{surplus}} = +0.082$, $+0.077$, $+0.058$ nats at rungs 1, 2, and 3 respectively, a mild decay but consistently positive, in contrast to the 1st-tier rebate which collapses and inverts by rung~3.
\end{quote}

Where Turkish's Agglutinative Compounding Effect (\S\ref{sec:results-naming}) describes the functional form of the rebate over $\rho \cdot H$ for surface-recoverable morphologies, the Templatic Lexical Surplus describes the decomposition of the rebate for morphologies whose grammatical information sits partly behind a lexical filter. The two named findings are complementary: they operate on different axes of the framework (functional form vs.\ decomposition) and are visible in the two typological poles that most differ from the English baseline of mainstream tokenization research.

A caveat: the 2nd-tier model has $\sim 1.3$\,M more parameters than the 1st-tier model (the added $20\mathrm{k} \times 64$ root embedding plus tied output), so the 0.082-nat gain is an upper bound on the attribution to the root stream alone. That said, the 1st-tier model at 3.3\,M parameters is already heavily over-parameterised relative to 5\,M training tokens, and there is no capacity shortage an additional 1.3\,M parameters would fill without the root-stream signal. The gain is therefore plausibly attributable in substantial part to the lexicon-recoverable axis becoming explicit at input; a parameter-matched ablation is scoped in \S\ref{sec:conclusion}.

This result validates the \S\ref{sec:discussion-limitations} decomposition of $\rho \cdot H$ into surface-recoverable and lexicon-recoverable components, giving Arabic back credit our surface-only $\rho$ refuses it. The 2nd-tier experiment is a language-specific tool rather than a universal extension: Turkish is already surface-recoverable, and English's gain would be smaller still.

\paragraph{A caveat on what the 2nd tier does \emph{not} encode.}
The 2nd-tier experiment exposes pure root identity to the model but does not expose the \emph{semantic function} of each wazn: the categorical interpretation a native reader derives from the pattern itself. Arabic's awzān are not neutral morphosyntactic scaffolding: the pattern wazn al-f\=a\textsuperscript{c}il (فاعِل) marks an agent noun, wazn al-maf\textsuperscript{c}\=ul (مَفْعُول) marks a patient, wazn al-\=alah (مِفْعَال, مِفْعَلَة) marks an instrument, the masdar patterns mark verbal nouns, the verbal Form X (اسْتَفْعَل) typically carries seeking or considering semantics, Form II (فَعَّل) intensification or causativity, and so on. These are first-class semantic classes that the pattern deterministically encodes for a human reader but which the current engine lumps into coarse templatic categories (\texttt{NOM\_DERIVED}, verbal Form I--X) without exposing the category's semantic function to the model. The 2nd-tier surplus reported here therefore understates the full Templatic Lexical Surplus: it counts only the rebate attributable to explicit root identity, not the rebate the wazn's semantic function would additionally contribute if exposed. A 3rd-tier extension that exposes wazn semantic class as a fourth input stream (a natural continuation of the decomposition) is scoped as follow-up in \S\ref{sec:conclusion}.

\FloatBarrier
\subsection{Compositional generalisation set for Arabic}
\label{sec:results-comp-gen}

A question Arabic can address that the other languages cannot is whether morph-aware tokenisation induces \emph{compositional learning} of the root-pattern system. A model that has learned the morphology should generalise to (root, template) combinations it never saw in training, provided each root and each template are individually attested. To support this evaluation we constructed a held-out (root, template) bucket on the Arabic test corpus: of 235{,}054 test-set word instances, 92.1\,\% have their pair attested in training (SEEN\_PAIR), 2.0\,\% combine a seen root with a seen template in an unseen combination (HELD\_OUT), and 5.8\,\% involve a root or template novel to training (NOVEL\_ATOM). Per-model loss attribution on this bucket requires word-to-token position tracking that differs across tokenisers and is scoped as follow-up in \S\ref{sec:conclusion}. A clean result would be that morph regimes show a \emph{larger} advantage over baseline on HELD\_OUT than on SEEN\_PAIR, evidence of compositional learning rather than memorised surface--bundle co-occurrences.

\subsection{Summary}
\label{sec:results-summary}

Phase~1 confirmed the framework's ordinal prediction and exposed the magnitude puzzle that motivated naming the Agglutinative Compounding Effect. The three-rung Chinchilla-ratio ladder sharpened the verdict. Turkish's morph advantage decays from $+0.770$ to $+0.535$ nats across rungs but stays substantially positive at rung~3, where the identical morph-vocab-cap confound drives English and Arabic below zero: the surviving Turkish rebate is the strongest evidence that at the high-$\rho \cdot H$ pole the morph mechanism is a real property of the language--tokenizer pairing (Interpretation~A), even with magnitude decay. For Arabic and English, the rung-3 and rung-2 inversions remain consistent with either Interpretation~B or the vocab-cap confound; a parameter-matched follow-up is scoped in \S\ref{sec:conclusion}. The two-tier Arabic extension of \S\ref{sec:results-tier2} adds a complementary finding on a different axis: an explicit root-embedding stream delivers an additional 0.082, 0.077, 0.058 nats at rungs~1--3, validating the decomposition of $\rho \cdot H$ into surface- and lexicon-recoverable components.

Together, the two named findings re-shape the central claim. The framework's $\rho \cdot H$ axis predicts not only cross-linguistic ordering at every rung (including sign-flipped regimes) but also the rate at which each language's rebate compresses as its baseline approaches its entropy floor. The \textbf{Agglutinative Compounding Effect} names the functional-form behaviour visible across the ladder; the \textbf{Templatic Lexical Surplus} names the complementary lexicon-axis decomposition. Together they map the boundary of what a surface-only tokenizer can exploit and what a lexicon-aware one can reach beyond.
